\documentclass[11pt]{amsart}

\usepackage[T1]{fontenc}
\usepackage[utf8]{inputenc}
\usepackage{amsmath,amssymb,amsthm}
\usepackage{hyperref}

\newcommand{\LL}{\Lambda}
\newcommand{\OO}{\mathbb{O}}
\newcommand{\ZZ}{\mathbb{Z}}
\newcommand{\RR}{\mathbb{R}}

\newtheorem*{observation}{Observation}
\newtheorem{example}{Example}[section]

\title[A companion note on the $\sigma$-twist]%
{A companion note\\
\large Kirmse--Coxeter and the $\sigma$-twist, by concrete example}
\author{Jens K\"oplinger}
\date{April 2026}

\begin{document}
\maketitle

\begin{abstract}
This note accompanies the paper ``An order on the Leech lattice from
transposition-twisted octonions.''  Its purpose is purely expository:
to walk through the construction step by step, pairing every abstract
definition with a concrete hand-checkable example, and to do so first
at the~$E_8$ level (where Coxeter's correction of Kirmse provides the
historical precedent) and then at the Leech level (where the
$\sigma$-twist plays the analogous role).  No new results are
proved.  Readers who are comfortable with the main paper will find
nothing new here; readers who want to see every piece verified by
hand will find the explicit numerics they may have missed.
\end{abstract}

% ===================================================================
\section{Preliminaries}
% ===================================================================

This note follows the conventions of the main paper.  We work with
the real octonion algebra~$\OO$ on $\RR^8$ with basis
$\{e_0, e_1, \ldots, e_7\}$, and with Wilson's construction of the
Leech lattice inside the triple octonion space
$\RR^{24} = \OO_1 \oplus \OO_2 \oplus \OO_3$.

The guiding principle throughout is: \emph{distinguish abstract
structure from the specific lattice to which it is applied}.  The
same ambient algebra can host many lattices; the same lattice can
be tested against many bilinear products on the ambient space.
``Closure of this lattice under this product'' is a question about
both, not just one.

% ===================================================================
\section{The $E_8$ lattice}\label{sec:e8}
% ===================================================================

\subsection*{Abstract}

The lattice~$E_8$ is the unique even unimodular lattice in~$\RR^8$
(up to isomorphism).  ``Even'' means every vector~$v \in E_8$ has
integer squared norm $\|v\|^2 \in 2\ZZ$.  ``Unimodular'' means
$E_8$ equals its own dual lattice (the set of vectors having
integer inner product with every element of~$E_8$).  It contains
$240$ vectors of minimal squared norm~$2$, the \emph{roots}.

\subsection*{Concrete}

Following Wilson, take for $L$ the integer span of the following
240 octonions:
\begin{itemize}
\item \textbf{112 ``type-1'' roots} of the form
  $\pm e_a \pm e_b$ with $0 \le a < b \le 7$;
\item \textbf{128 ``type-2'' roots} of the form
  $\tfrac{1}{2}(\varepsilon_0 e_0 + \varepsilon_1 e_1 + \cdots
  + \varepsilon_7 e_7)$ where each
  $\varepsilon_i \in \{+1, -1\}$ and the number of~$-1$s is odd.
\end{itemize}

Equivalently,
\[
  L = \{v \in \ZZ^8 : \sum_i v_i \text{ is even}\}
   \;\cup\;
      \{v \in (\ZZ + \tfrac{1}{2})^8 : \sum_i v_i \text{ is odd}\}.
\]

\begin{example}
The vector $(1, 1, 0, 0, 0, 0, 0, 0) = e_0 + e_1$ is a type-1 root:
integer coordinates, sum~$2$.  The vector
$\tfrac{1}{2}(-1, 1, 1, 1, 1, 1, 1, 1)$ is a type-2 root:
half-integer coordinates, one $-1$ (odd count), coordinate sum $3$
(odd).  Both lie in~$L$; both have squared norm~$2$.
\end{example}

\begin{example}
Distinguished element in~$L$ to be used throughout:
\[
  s \;=\; \tfrac{1}{2}(-e_0 + e_1 + e_2 + e_3 + e_4 + e_5 + e_6 + e_7),
\]
of squared norm~$2$.  In coordinates $s = (-\tfrac{1}{2},
\tfrac{1}{2}, \tfrac{1}{2}, \tfrac{1}{2}, \tfrac{1}{2}, \tfrac{1}{2},
\tfrac{1}{2}, \tfrac{1}{2})$: half-integer, one~$-1/2$ (odd count),
sum~$3$ (odd).  Confirmed in~$L$.
\end{example}

% ===================================================================
\section{The octonion algebra}\label{sec:oct}
% ===================================================================

\subsection*{Abstract}

The octonions~$\OO$ are the unique real \emph{normed division
algebra} of dimension~$8$: a real $8$-dimensional vector space with
a bilinear product~$x \cdot y$ and a positive-definite quadratic
norm satisfying $\|x \cdot y\| = \|x\| \cdot \|y\|$.  The algebra is
non-commutative and non-associative (but alternative).

\subsection*{Concrete}

Fix the basis $\{e_0, e_1, \ldots, e_7\}$ with $e_0$ the
multiplicative identity.  The multiplication on the imaginary units
is determined by the seven \emph{Fano triples}
\begin{equation}\label{eq:fano}
  (1,2,4),\; (2,3,5),\; (3,4,6),\; (4,5,7),\; (5,6,1),\; (6,7,2),\;
  (7,1,3),
\end{equation}
read cyclically: for $(a, b, c)$ above, $e_a \cdot e_b = +e_c$,
$e_b \cdot e_c = +e_a$, $e_c \cdot e_a = +e_b$, with opposite signs
for the reversed products.  Every imaginary unit squares to~$-e_0$.

\begin{example}
From the triple~$(1,2,4)$: $\,e_1 \cdot e_2 = +e_4$ and
$e_2 \cdot e_1 = -e_4$.  From~$(7,1,3)$: $\,e_1 \cdot e_3 = +e_7$.
From~$(6,7,2)$: $\,e_7 \cdot e_2 = +e_6$, so
$e_2 \cdot e_7 = -e_6$.
\end{example}

\subsection*{Closure of $L$ under the octonion product}

Under the Fano convention \eqref{eq:fano}, the lattice~$L$ is closed
under octonion multiplication: for every $u, v \in L$, the product
$u \cdot v \in L$.  This is a classical result of Coxeter~(1946),
verified in the main paper by computing all $8 \times 8 = 64$
products of a $\ZZ$-basis of~$L$ in exact arithmetic and checking
that each lies again in~$L$.

\begin{example}
A sample product.  From $s \cdot s$ with $s$ as above, using the
identity $(\sum_{k=1}^{7} e_k)^2 = -7 e_0$:
\[
  s \cdot s \;=\;
    -\tfrac{3}{2}\, e_0 - \tfrac{1}{2}(e_1 + e_2 + e_3 + e_4 + e_5
    + e_6 + e_7).
\]
Its coordinates are all half-integer; the coordinate sum is
$-\tfrac{3}{2} - 7 \cdot \tfrac{1}{2} = -5$, odd.  Hence
$s \cdot s \in L$.
\end{example}

% ===================================================================
\section{The Kirmse--Coxeter precedent}\label{sec:kirmse}
% ===================================================================

\subsection*{Historical remark}

Kirmse (1924) proposed a specific set of 240 integral octonions
together with a multiplication table.  He asserted the set was
closed under the multiplication.  Coxeter (1946) noticed a subtle
error: Kirmse's multiplication table, as written, did not produce
closure.  Coxeter showed that if one permutes two basis indices in
the Fano-plane labelling, one obtains a (different) multiplication
table under which the 240-element set \emph{is} closed.  There are
seven such Coxeter-corrected conventions, one for each pair of
imaginary indices; each gives one of the seven maximal orders of
the integral octonions discussed in Conway--Smith.

The Fano triples \eqref{eq:fano} used in this note and in the main
paper are one of Coxeter's seven corrected conventions.  Under this
convention $L$ closes, as stated above.  We do not reproduce
Kirmse's original table here; the point of recalling the story is
the structural pattern, not the historical arithmetic:

\begin{observation}
A convention of octonion multiplication fixes how basis products
are assigned.  A different convention, obtained by transposing two
basis indices, assigns products differently.  For some lattices, the
two conventions are interchangeable (the lattice is closed under
both); for others, only one convention closes.  Which lattices
behave which way is precisely the question that Kirmse got wrong
and Coxeter got right.
\end{observation}

\subsection*{A concrete lattice where the story does matter}

For the lattice~$L$ itself, this distinction is invisible: $L$ is
symmetric under every coordinate permutation, so transposing two
basis indices produces an algebra under which $L$ also closes.  To
see the Kirmse--Coxeter phenomenon concretely at the $E_8$ level we
need a lattice that is \emph{not} symmetric under such a swap.  One
appears naturally in Wilson's Leech construction: the sublattice
\[
  Ls \;=\; \{\ell \cdot s : \ell \in L\},
\]
the image of $L$ under right-multiplication by the element~$s$
above.  This is again a full-rank sublattice of~$L$, of index~$16$
(see the main paper, Section~2.3).  It is a lattice, so of course it
is closed under addition; the question is whether it is closed under
octonion multiplication.

\begin{example}[Standard product fails on~$Ls$]\label{ex:Ls-fails}
Take
\[
  a = (-1,\, 0,\, 0,\, 0,\, 1,\, 0,\, 1,\, 1), \qquad
  b = (-1,\, 1,\, 0,\, 0,\, 0,\, 1,\, 0,\, 1).
\]
Both $a$ and $b$ are in $Ls$ (they appear as basis vectors of~$Ls$
produced in the main paper's exact-arithmetic verification).  Their
octonion product, computed under the Fano convention
\eqref{eq:fano}, is
\[
  a \cdot b = (0,\, -2,\, 2,\, 1,\, -2,\, -1,\, -1,\, -1).
\]
This vector has integer coordinates and lies in~$L$ --- but it is
\emph{not} an integer combination of a $\ZZ$-basis of~$Ls$.  The
verification is carried out with exact rational arithmetic in the
script
\href{../python_project/src/symbolic_proof_checks.py}{\texttt{symbolic\_proof\_checks.py}}.

Consequently, $Ls \cdot Ls \not\subseteq Ls$.
\end{example}

This is the structural echo of Kirmse's situation at the sublattice
level: a specific sublattice, under a specific octonion product,
fails to close.

% ===================================================================
\section{The Coxeter transposition, abstractly}\label{sec:swap}
% ===================================================================

\subsection*{Abstract}

Fix a transposition $\sigma = (s\; t)$ on the imaginary basis
indices $\{1, \ldots, 7\}$ (for example, $\sigma = (1\;2)$).
Extend~$\sigma$ linearly to a map $\RR^8 \to \RR^8$ that fixes
$e_0$ and swaps $e_s \leftrightarrow e_t$.  This~$\sigma$ is an
\emph{isometry} of~$\RR^8$: it preserves the Euclidean inner
product.  It is also an involution: $\sigma^2 = \mathrm{id}$.

Starting from the standard product~$\cdot$, the map~$\sigma$ induces
a \emph{$\sigma$-twisted product} on the same $\RR^8$:
\begin{equation}\label{eq:twist}
  x \cdot_\sigma y \;=\; \sigma\bigl(\sigma(x) \cdot \sigma(y)\bigr).
\end{equation}
The pair $(\RR^8, \cdot_\sigma)$ is again an octonion algebra,
abstractly isomorphic to $(\RR^8, \cdot)$ via~$\sigma$ itself.  But
it is a \emph{different bilinear product} on the same $\RR^8$: it
assigns different specific products to pairs of basis vectors.

\subsection*{Concrete}

Take $\sigma = (1\; 2)$, which swaps $e_1$ and~$e_2$ (and fixes
$e_0, e_3, e_4, e_5, e_6, e_7$).  Then
\begin{align*}
  e_1 \cdot_\sigma e_2
    &= \sigma(\sigma(e_1) \cdot \sigma(e_2))
     = \sigma(e_2 \cdot e_1)
     = \sigma(-e_4)
     = -e_4,
\end{align*}
whereas $e_1 \cdot e_2 = +e_4$ under the standard product.  Similarly
one checks $e_1 \cdot_\sigma e_4 = +e_2$ (vs.\ $-e_2$ under
$\cdot$), and so on: the two products disagree on some pairs of
basis elements.

\subsection*{What is twisted and what is invariant}

\begin{itemize}
\item \textbf{Twisted:} the multiplication table.  Under
  $\cdot_\sigma$, the Fano triples \eqref{eq:fano} are replaced by
  $(\sigma(a), \sigma(b), \sigma(c))$ for each original triple.
\item \textbf{Invariant under~$\sigma$ (as a linear map):}
  the underlying vector space~$\RR^8$, the Euclidean inner product,
  the decomposition into integer and half-integer vectors, and the
  lattice~$L$ itself ($\sigma(L) = L$ since $\sigma$ is a coordinate
  permutation).
\item \textbf{Not invariant under~$\sigma$:} specific sublattices
  of~$L$ that depend on a distinguished element, such as~$Ls$.
  Since $Ls$ is defined as $\{\ell \cdot s : \ell \in L\}$, applying
  $\sigma$ gives
  \[
    \sigma(Ls) = \{\sigma(\ell) \cdot_{\sigma} \sigma(s)
                    : \ell \in L\}
    \;\ne\; Ls
  \]
  in general.  The main paper verifies $\sigma(Ls) \ne Ls$
  directly: for $\sigma = (1\;2)$, the vector
  $(-1, 0, 1, 0, 0, 1, 0, 1)$ lies in~$\sigma(Ls)$ but not
  in~$Ls$.
\end{itemize}

The last two bullets together say: \emph{$\sigma$ is an isometry of
$\RR^8$ that fixes~$L$ but moves~$Ls$ to a distinct sublattice.}

% ===================================================================
\section{The twist at the $E_8$ level changes the story}
% ===================================================================

The sublattice $\sigma(Ls)$ is the image of~$Ls$ under the
transposition.  It has the same rank as~$Ls$, the same index in~$L$,
and is isometric to~$Ls$.  But its elements are different vectors.
The question of closure under the standard octonion product is now
posed afresh.

\begin{example}[Standard product succeeds on~$\sigma(Ls)$]
\label{ex:sigmaLs-succeeds}
A $\ZZ$-basis of~$\sigma(Ls)$ is obtained by applying~$\sigma$
(coordinate swap of positions $1, 2$) to each basis vector of~$Ls$.
Using the basis generated in the main paper's verification, all
$8 \times 8 = 64$ products of basis vectors of~$\sigma(Ls)$, under
the standard octonion product~$\cdot$, lie again in~$\sigma(Ls)$.
The verification is in
\href{../python_project/src/symbolic_proof_checks.py}{\texttt{symbolic\_proof\_checks.py}}
as Lemma~D.

Hence $\sigma(Ls) \cdot \sigma(Ls) \subseteq \sigma(Ls)$.
\end{example}

\subsection*{Two equivalent readings}

Because $\sigma$ is an isometry and an algebra isomorphism
$(\RR^8, \cdot) \to (\RR^8, \cdot_\sigma)$, one can restate
Example~\ref{ex:sigmaLs-succeeds} in an equivalent form on the
original lattice~$Ls$ with the twisted product:
\[
  Ls \cdot_\sigma Ls \;\subseteq\; Ls.
\]
(Apply $\sigma$ to both sides of the inclusion
$\sigma(Ls) \cdot \sigma(Ls) \subseteq \sigma(Ls)$ and use
$\sigma^2 = \mathrm{id}$ together with
\eqref{eq:twist}.)  The two forms are tautologically equivalent.

Reading both against Example~\ref{ex:Ls-fails}:
\[
  \underbrace{Ls \cdot Ls \not\subseteq Ls}_{\text{standard product
  fails on $Ls$}}
  \qquad\text{versus}\qquad
  \underbrace{Ls \cdot_\sigma Ls \subseteq Ls}_{\text{twisted product
  succeeds on $Ls$}}.
\]
Equivalently:
\[
  \underbrace{Ls \cdot Ls \not\subseteq Ls}_{\text{$Ls$ fails}}
  \qquad\text{versus}\qquad
  \underbrace{\sigma(Ls) \cdot \sigma(Ls) \subseteq \sigma(Ls)}
    _{\text{$\sigma(Ls)$ succeeds}}.
\]

\subsection*{Interpretation}

Two different bilinear products on~$\RR^8$ (namely~$\cdot$
and~$\cdot_\sigma$) are isomorphic as abstract algebras.  Every
theorem phrased in the language of the abstract algebra carries
over.  But ``does lattice~$X$ close under product~$\star$?'' is not
a theorem phrased purely in the abstract algebra: it also names a
specific subset $X \subset \RR^8$, and the answer depends on how the
isomorphism relates $X$ to itself.

If $\sigma(X) = X$, the answer is the same for both products.  If
$\sigma(X) \ne X$, the answer can flip: $X$ may fail to close
under~$\cdot$ while $\sigma(X)$ closes under~$\cdot$; equivalently,
$X$ closes under $\cdot_\sigma$ while~$\sigma(X)$ fails under
$\cdot_\sigma$.  The lattice~$L$ itself is $\sigma$-invariant and
therefore closes under both products.  The sublattice~$Ls$ is
\emph{not} $\sigma$-invariant, and this is exactly what opens the
gap between the two closure statements.

This is the structural lesson of Kirmse--Coxeter in isolation.  The
same lesson, applied one dimension higher, is the content of the
main paper.

% ===================================================================
\section{Wilson's Leech lattice}\label{sec:wilson-leech}
% ===================================================================

\subsection*{Abstract}

The Leech lattice~$\LL$ is the unique even unimodular lattice of
rank~$24$ with no roots (i.e.\ no vectors of squared norm~$2$).  Its
minimal vectors have squared norm~$4$, and there are exactly
$196{,}560$ of them.

\subsection*{Concrete (Wilson's construction)}

Identify $\RR^{24}$ with $\OO_1 \oplus \OO_2 \oplus \OO_3$: three
copies of the octonions, labelled $\alpha = 1, 2, 3$.  Let
$\bar{s} = \tfrac{1}{2}(-e_0 - e_1 - e_2 - \cdots - e_7)$ (the
octonion conjugate of the negative of the earlier~$s$; it also lies
in~$L$).  Define two sublattices of~$L$:
\[
  Ls = \{\ell \cdot s : \ell \in L\}, \qquad
  L\bar{s} = \{\ell \cdot \bar{s} : \ell \in L\},
\]
each of index~$16$ in~$L$.  Wilson's Leech lattice is then
\begin{multline}\label{eq:leech}
  \LL \;=\; \bigl\{(x, y, z) \in L^3 \;:\\
    x + y,\; x + z,\; y + z \;\in\; L\bar{s}; \quad
    x + y + z \;\in\; Ls \bigr\}.
\end{multline}

The sublattice~$Ls$ from Section~\ref{sec:kirmse} is precisely the
``condition-(3) block sum'' lattice in Wilson's definition.

% ===================================================================
\section{The standard triple product on $\RR^{24}$}
% ===================================================================

\subsection*{Abstract}

The natural bilinear product on $\RR^{24} = \OO_1 \oplus \OO_2
\oplus \OO_3$ built from a single octonion product is a
\emph{triple product with $\ZZ_3$-symmetric cross-block routing}.
For $u = (x, y, z)$ and $v = (x', y', z')$ define
\[
  u \star v \;=\; (P_1, P_2, P_3),
\]
where
\begin{align*}
  P_1 &= x \cdot x' + z \cdot y' + y \cdot z', \\
  P_2 &= y \cdot y' + x \cdot z' + z \cdot x', \\
  P_3 &= z \cdot z' + y \cdot x' + x \cdot y'.
\end{align*}
Same-block products stay in the block; cross-block products
$O_\alpha \times O_\beta$ land in $O_\gamma$ with
$\{\alpha, \beta, \gamma\} = \{1, 2, 3\}$.  The three blocks play
the same role under cyclic rotation $1 \to 2 \to 3 \to 1$.

\subsection*{Why condition~(3) is the obstruction}

A direct computation from the definitions gives
\[
  P_1 + P_2 + P_3 \;=\; (x + y + z) \cdot (x' + y' + z').
\]
Thus ``$u \star v$ satisfies Wilson's condition~(3)'' amounts to:
``the product of two elements of~$Ls$ lies in~$Ls$''.

Using $u, v \in \LL$ means $x + y + z \in Ls$ and $x' + y' + z'
\in Ls$, so the question reduces exactly to whether $Ls \cdot Ls
\subseteq Ls$.  By Example~\ref{ex:Ls-fails}, it does not.  Hence:

\begin{example}[Standard triple product fails on~$\LL$]
\label{ex:triple-fails}
There exist $u, v \in \LL$ for which $u \star v \notin \LL$: the
block sum of $u \star v$ need not lie in~$Ls$, failing Wilson's
condition~(3).  The concrete obstruction is precisely the pair
$(a, b)$ of Example~\ref{ex:Ls-fails}: any $u, v \in \LL$ whose
triple sums equal~$a, b$ respectively produce a $u \star v$ whose
block sum equals~$a \cdot b \notin Ls$.
\end{example}

% ===================================================================
\section{The $\sigma$-twisted triple product}
% ===================================================================

\subsection*{Abstract}

Apply~$\sigma$ componentwise to each of the three octonion blocks.
This yields a linear map on~$\RR^{24}$, still denoted~$\sigma$, that
is an isometry and satisfies $\sigma^2 = \mathrm{id}$.  Replace the
octonion product in each of the nine block-pair products of~$\star$
by the $\sigma$-twisted product $\cdot_\sigma$ from
\eqref{eq:twist}: this produces a new bilinear product
$\star_\sigma$ on~$\RR^{24}$.  By the same formal manipulation as
before,
\[
  u \star_\sigma v \;=\; \sigma\bigl(\sigma(u) \star \sigma(v)\bigr),
\]
where on the right $\star$ is the standard (untwisted) triple
product.

\subsection*{Concrete}

Choose $\sigma = (1\; 2)$, acting on imaginary indices of each
octonion block.  Then for $u, v \in \LL$,
\[
  (\text{block sum of } u \star_\sigma v)
  \;=\;
  \sigma\bigl((\sigma(x) + \sigma(y) + \sigma(z)) \cdot
              (\sigma(x') + \sigma(y') + \sigma(z'))\bigr).
\]
Inside the outer~$\sigma$, each triple sum lies in~$\sigma(Ls)$
(since each original triple sum lies in~$Ls$, by Wilson's
condition~(3) on~$u, v$).  By
Example~\ref{ex:sigmaLs-succeeds}, $\sigma(Ls) \cdot \sigma(Ls)
\subseteq \sigma(Ls)$.  Applying the outer~$\sigma$ brings the
result into $\sigma(\sigma(Ls)) = Ls$.

So the block sum of $u \star_\sigma v$ lies in~$Ls$: Wilson's
condition~(3) is restored.  Conditions~(1) and~(2) go through by
analogous arguments (see the main paper, Section~4).  The net effect
is:

\begin{example}[$\sigma$-twisted triple product succeeds on~$\LL$]
For all $u, v \in \LL$, the product $u \star_\sigma v$ lies
in~$\LL$.  This is Theorem~3.1 of the main paper.
\end{example}

\subsection*{What is twisted and what is invariant (at the Leech
level)}

\begin{itemize}
\item \textbf{Twisted:} the octonion multiplication used inside
  each of the nine block-pair products.  Equivalently, the Fano
  triples~\eqref{eq:fano} are replaced block by block by their
  $\sigma$-images.
\item \textbf{Invariant under~$\sigma$:} the ambient space
  $\RR^{24}$, the three-block decomposition into~$\OO_1, \OO_2,
  \OO_3$, the $\ZZ_3$ cross-block routing rule, and the
  sub-lattice~$L$ of each block ($\sigma(L) = L$).
\item \textbf{Not invariant under~$\sigma$:} the sub-lattice~$Ls$
  of each block ($\sigma(Ls) \ne Ls$), and through it the Leech
  lattice itself ($\sigma(\LL) \ne \LL$ in general, since~$\LL$ is
  defined by a condition involving~$Ls$).
\end{itemize}

The comparison with the $E_8$ level is exact.  One level down,
$\sigma$ fixes~$L$ but not~$Ls$, and this is what allows the twist
to change the closure question for~$Ls$.  One level up, $\sigma$
fixes~$\OO_1 \oplus \OO_2 \oplus \OO_3$ but not~$\LL$, and this is
what allows the twist to change the closure question for~$\LL$.

% ===================================================================
\section{Could the $\sigma$-twist undo Coxeter's fix?}
% ===================================================================

A natural concern:\ Coxeter's correction~$\kappa$ of Kirmse swaps two
Fano triples in the multiplication table, producing a convention
under which the $E_8$ lattice~$L$ closes.  Our~$\sigma$ is also a
transposition.  Could a poor choice of~$\sigma$ effectively reverse
Coxeter's swap and return us to Kirmse's 1924 regime, where~$L$
fails to close?

The answer is no, for a reason that is worth surfacing because it
also explains why~$\sigma$ can help at all.  Coxeter's~$\kappa$ and
our~$\sigma$ act at different structural levels:

\begin{itemize}
\item Coxeter's~$\kappa$ is a transposition of \emph{Fano
  triples}.\ It modifies the octonion multiplication table globally:\
  the products of \emph{all} basis vectors whose Fano lines are
  touched are rewritten.
\item Our~$\sigma$ is a transposition of \emph{coordinate axes}.\ It
  is a linear involution of $\RR^8$ that extends to the multiplication
  by conjugation, $a \cdot_\sigma b = \sigma(\sigma(a) \cdot
  \sigma(b))$.
\end{itemize}

Most coordinate-level transpositions do \emph{not} correspond to
Fano-triple permutations; for instance, with our conventions
$\sigma = (1\;2)$ does not send the triple $(2,3,5)$ to any of the
other six listed in~\eqref{eq:fano}, so the Fano plane's line
structure is not preserved by~$\sigma$ at all.\ In particular,
$\sigma$ cannot be $\kappa^{-1}$\ because~$\kappa$ \emph{is} a
rearrangement of lines while~$\sigma$ is not.

Beyond this structural separation, there is a stronger and more
directly relevant guarantee.  The sub-lattice~$L$ is invariant
under~$\sigma$:
\[
  \sigma(L) = L
  \qquad \text{(Lemma A of the main paper).}
\]
Coxeter's correction is a statement about~$L$:\ under his corrected
convention, $L \cdot L \subseteq L$.\ Because~$\sigma$ fixes~$L$, the
$\sigma$-twisted product inherits this closure directly:
\[
  L \cdot_\sigma L
  \;=\; \sigma\bigl(\sigma(L) \cdot \sigma(L)\bigr)
  \;=\; \sigma(L \cdot L)
  \;\subseteq\; \sigma(L)
  \;=\; L.
\]
So regardless of which coordinate transposition~$\sigma$ we choose
(so long as $\sigma(L) = L$, which holds for any permutation of the
imaginary axes that preserves the $D_8^+$ structure), Coxeter's $E_8$
closure is automatically preserved.  Kirmse's 1924 regime, which
failed at the level of~$L$ itself, cannot be reached from Coxeter's
convention by a twist that fixes~$L$.

What the $\sigma$-twist \emph{does} change is how the product sees
sub-lattices of~$L$ that~$\sigma$ does not fix --- in our case,
$Ls$.\ Here~$\sigma(Ls) \ne Ls$ (explicit witness in item~1 of the
next section), and this is precisely what allows the twist to
resolve Wilson's
condition~(3) without touching conditions~(1) or~(2).  The two
corrections --- Coxeter's~$\kappa$ at the Fano-triple level, and
our~$\sigma$ at the coordinate level --- are therefore acting on
different strata of the same construction:\ Coxeter's fixes the
$E_8$ stratum once and for all; ours redirects a strictly
finer-grained sub-lattice question that Coxeter's convention, while
necessary, does not by itself settle.

In short: Coxeter's~$\kappa$ is not a linear map on~$\RR^8$ at all
(it is a rearrangement of the multiplication table); our~$\sigma$
is a linear map on~$\RR^8$ that happens to fix~$L$ on the nose.
There is no route along the $\sigma$-family back to Kirmse's regime.

% ===================================================================
\section{Summary of concrete data}
% ===================================================================

\noindent Three empirical facts are verified by exact rational
arithmetic in
\href{../python_project/src/symbolic_proof_checks.py}
     {\texttt{symbolic\_proof\_checks.py}}:
\begin{enumerate}
\item $\sigma(Ls) \ne Ls$ (explicit witness:
  $(-1, 0, 1, 0, 0, 1, 0, 1) \in \sigma(Ls) \setminus Ls$).
\item $Ls \cdot Ls \not\subseteq Ls$ under the standard octonion
  product (explicit witness~$a, b$ in
  Example~\ref{ex:Ls-fails} with $a \cdot b \notin Ls$).
\item $\sigma(Ls) \cdot \sigma(Ls) \subseteq \sigma(Ls)$ under the
  standard octonion product (all 64 basis products land in
  $\sigma(Ls)$).
\end{enumerate}

\noindent Two equivalent phrasings of the main result are:
\begin{itemize}
\item ``$\LL$ is closed under the $\sigma$-twisted triple product
  $\star_\sigma$.''
\item ``$\sigma(\LL)$ is closed under the standard triple product
  $\star$.''
\end{itemize}
Both say the same thing via the algebra isomorphism
$\sigma: (\RR^{24}, \star) \to (\RR^{24}, \star_\sigma)$; neither
says that $\LL$ is closed under~$\star$, which is false.

% ===================================================================

\end{document}
